% ALl of your information should be placed here

%
\author{Der Autor} %z.b. Max Musterman
\def\MatrNummer{5} % example for a matrikelnumber
\def\Studiengang{IT-Security}  % bsp IT
\def\Email{der.autor@example.com}
\def\Addresse{Adresse des Autors} % Hauptstr. 1.\\ 100000 Berlin % DIe Unterbrechung sollte man machen muss man aber nicht


% Spezial Variable für bachlor Template
\def\Abschlussbezeinung{Bachelor of Sience IT-Security}


% Pruefungsrelevante Fragen
\def\Pruefer{Gutachter}  % Prof Max Musterman
\def\ZweitGutachter{Zweitgutachter} % Nur Relevant bei Bachlorarbeit
\def\ArtArbeit{Bachelorarbeit} % Bachlor Thesis, Hausarbeit etc.
\def\Modul{} % Dein Modul in dem du das ganze schreibt
\def\Semester{SS2026} %  Als beispielWS2023/2024
\title{Die Grenzen von cachebasierten Covert-Channels im Kontext von Transient-Execution: ein technischer Überblick} %your title


% usually unchanged things
\date{1.6.2026}


\def\Firma{Firma} %z.B. Telekom
\def\Hochschulname{Leibniz-Fachhochschule}

% Toggels
\def\Sperrvermerk{0} % definiert ob ein Sperrvermerkt benutzt werden
\def\Acronyms{1} % definiert ob das Abkürzungsverzeichnis benutzt wird
\def\BachlorThesis{1} % definiert ob es eine Bachlor Thesis ist oder nicht
\def\digitalSig{1} % definiert ob eine digitale Signatur verwendet werden soll (ergo mit adobe signieren ) ansonsten wird einfach ein bild eingebunden
\def\parenthesisCite{0} % definiert ob alle \cite automatisch durch \parencite ersetzt werden sollen, sodass aus Autor Jahr -> (Autor, Jahr) wird

% Optionales bearbeiten von Variablen
\def\acrotitle{Abkürzungsverzeichnis}

% die Pfade für die BIB ressource und diese Datei müssen lediglich angepasst werden in der main.tex ansonsten nichts.
